\documentclass[12pt, a4paper]{article}
\usepackage[utf8]{inputenc}
\usepackage[T1]{fontenc}
\usepackage[ngerman]{babel}
\usepackage{amsmath, amssymb, amsthm, mathtools}
\usepackage{graphicx}
\usepackage{booktabs}
\usepackage{array}
\usepackage{geometry}
\usepackage{setspace}
\usepackage{pgfplots}
\pgfplotsset{compat=1.18}
\usepackage{tikz}
\usetikzlibrary{arrows.meta, positioning, shapes.geometric, calc}
\usepackage{algorithm}
\usepackage{algpseudocode}
\usepackage{enumitem}
\setlist[itemize,1]{label=--}
\usepackage{xcolor}
\usepackage{caption}
\usepackage{subcaption}
\usepackage{listings}
\usepackage{lmodern}
\usepackage{microtype}
\usepackage{xcolor}
\usepackage[hidelinks]{hyperref}
\usepackage{cleveref}
\geometry{margin=2.5cm}

\hypersetup{colorlinks=true,linkcolor=blue!60!black,urlcolor=blue!60!black,citecolor=blue!60!black}

% ---- cleveref: deutsche Kapitel-Verweise ----
\crefname{section}{Kapitel}{Kapitel}
\Crefname{section}{Kapitel}{Kapitel}


% ---- Theorem-Umgebungen ----
\newtheorem{theorem}{Theorem}[section]
\newtheorem{lemma}[theorem]{Lemma}
\newtheorem{proposition}[theorem]{Proposition}
\newtheorem{corollary}[theorem]{Korollar}
\newtheorem{definition}[theorem]{Definition}
\newtheorem{example}[theorem]{Beispiel}
\newtheorem{remark}[theorem]{Bemerkung}

% ---- Operatoren / Notation ----
\DeclareMathOperator{\E}{\mathbb{E}}
\DeclareMathOperator{\Var}{Var}
\DeclareMathOperator{\Cov}{Cov}
\DeclareMathOperator{\Prob}{\mathbb{P}}
\newcommand{\Geo}{\mathrm{Geo}}
\newcommand{\Exp}{\mathrm{Exp}}
\newcommand{\Pois}{\mathrm{Pois}}
\newcommand{\NBin}{\mathrm{NegBin}}
\newcommand{\Pareto}{\mathrm{Pareto}}

% ---- Listings (Python-artig) ----
\lstset{
  language=Python,
  basicstyle=\ttfamily\small,
  keywordstyle=\color{blue},
  commentstyle=\color{gray},
  stringstyle=\color{red!60!black},
  numbers=left,
  numberstyle=\tiny\color{gray},
  frame=single,
  breaklines=true,
  captionpos=b
}

% ============================================================
\begin{document}

\begin{titlepage}
  \centering
  \vspace*{2cm}
  {\LARGE\bfseries Das Archiv-Problem\\[0.4em]}
  {\large Stochastische Wartestrategien zur Minimierung von\\
          Archivverschiebungen mit Anwendungen in der\\
          Rechnerarchitektur und Speicherverwaltung\\[2.5em]}
  \hrule height 1pt
  \vspace{1.5em}
  {\large\itshape Eine umfassende mathematische Analyse von Stephan Epp\\[1em]}
  \hrule height 0.5pt
  \vfill
  {\large 06. März 2026}
\end{titlepage}
\newpage

\tableofcontents
\newpage

% ============================================================
\section{Einleitung}\label{sec:intro}
% ============================================================
Das \emph{Archiv-Problem} beschreibt die Herausforderung, eine geordnete
Sammlung von Dokumenten -- das Archiv -- mit minimaler Anzahl von
Positionsverschiebungen zu pflegen, wenn die zukünftige Ankunftsreihenfolge
der einzuordnenden Dokumente \emph{a priori} unbekannt ist.
Anstatt die unbekannte Zukunft vorherzusagen, schlägt diese Arbeit eine
\emph{rückwärtsgewandte stochastische Wartestrategie} vor: Nach dem Eintreffen
eines Dokuments wird eine zufällige Wartezeit abgewartet; trifft in dieser Zeit
ein weiteres Dokument ein, beginnt die Wartezeit neu. Erst wenn eine
vollständige Warteperiode ohne neues Dokument verstreicht, wird die akkumulierte
Menge optimal einsortiert. Wir untersuchen diese Strategie unter Exponential-,
Geometrischer, Poisson-, Negativbinomial- und Pareto-Verteilung, leiten
geschlossene Ausdrücke für erwartete Wartezeiten und Verschiebungskosten her
und zeigen, welche Verteilung welches Systemziel optimiert. Abschließend wird
der direkte Bezug zur Cache-Hierarchie, zur Hauptspeicherverwaltung und zu
modernen Speichercontrollern in der Rechnerarchitektur hergestellt.

Physische und digitale Archive stehen vor einem grundlegenden
Optimierungsproblem: Eine geordnete Folge von Dokumenten soll so gepflegt
werden, dass der Suchaufwand minimiert wird. Bei einer alphabetischen oder
sachlichen Ordnung bedeutet das Einordnen eines neuen Dokuments oft, dass
bereits vorhandene Dokumente verschoben werden müssen – im schlimmsten Fall
alle. Bei $n$ Dokumenten im Archiv kann ein einziges neues Dokument bis zu
$\Theta(n)$ Verschiebungen erzwingen.

Das eigentliche Problem ergibt sich aus der \emph{Online-Natur} des Prozesses:
Neue Dokumente treffen sukzessive ein, und zum Zeitpunkt der Einordnung ist
die zukünftige Reihenfolge unbekannt. Würde man alle künftig eintreffenden
Dokumente kennen (\emph{Offline-Szenario}), könnte man eine nahezu
verschiebungsfreie Strategie planen (z.\,B. Batch-Sortierung). Da jedoch die
Zukunft unbekannt ist (\emph{Online-Szenario}), liegt es nahe, einen
\emph{Beobachtungshorizont aus der Vergangenheit} zu nutzen, um die aktuelle
Ankunftsrate zu schätzen, und dann eine zufällige Zeit zu warten, bevor man
einordnet. Trifft in dieser Wartezeit ein weiteres Dokument ein, wird die
Wartezeit zurückgesetzt – das System sammelt Dokumente, bis eine
``ruhige'' Periode eintritt.

Diese Idee ist in der Informatik unter verschiedenen Gewändern bekannt:
Ethernet-CSMA/CD-Backoff, adaptive Polling in Betriebssystemen, Batching in
Datenbankpuffern, sowie Interrupt-Coalescing in Netzwerkkarten. Das Archiv
dient hier als abstraktes Modell, das sich unmittelbar auf Speicherhierarchien
in der Rechnerarchitektur übertragen lässt.

\paragraph{Beiträge dieser Arbeit.}
\begin{itemize}[noitemsep]
  \item Formale Definition des Archiv-Problems als Online-Sortieraufgabe mit
        zufälliger Wartestrategie (\cref{sec:model}).
  \item Umfassende mathematische Analyse der Wartezeit-Verteilungen
        (Exponential, Geometrisch, Poisson, Negativbinomial, Pareto) mit
        geschlossenen Ausdrücken für Erwartungswert, Varianz und
        Verteilungsfunktion (\cref{sec:distributions}).
  \item Beweis der Optimalitätseigenschaften und Zieldifferenzierung zwischen
        den Verteilungen (\cref{sec:optimality}).
  \item Anwendung auf Cache-Hierarchien, TLB-Miss-Behandlung und
        Speichercontroller-Pufferung (\cref{sec:architecture}).
  \item Numerische Experimente und Simulation (\cref{sec:simulation}).
\end{itemize}

% ============================================================
\section{Problemformulierung und Modell}\label{sec:model}
% ============================================================

\subsection{Das Archiv als geordnete Datenstruktur}

\begin{definition}[Archiv]\label{def:archive}
Ein \emph{Archiv} $\mathcal{A}$ ist eine linear geordnete Liste von
$m$~Dokumenten $d_1 \prec d_2 \prec \cdots \prec d_m$ gemäß einer totalen
Ordnung $\prec$ (z.\,B. alphabetisch nach Schlüssel). Das \emph{Einfügen}
eines neuen Dokuments~$d$ an Position~$k$ in eine Liste der Länge~$m$
erfordert $m - k + 1$ \emph{Verschiebungen} aller Dokumente $d_k, d_{k+1},
\ldots, d_m$ um eine Position nach rechts.
\end{definition}

\begin{definition}[Verschiebungskosten]\label{def:cost}
Die \emph{Verschiebungskosten} für das Einfügen von Dokument $d$ in ein
Archiv der Länge~$m$ sind
\[
  C(d, \mathcal{A}) = \bigl|\{i : d_i \prec d\}\bigr| \;=\; \mathrm{rank}_{\mathcal{A}}(d)
  \quad\in\{0, 1, \ldots, m\}.
\]
Die Gesamtkosten für das sequentielle Einfügen einer Folge
$\sigma = (d^{(1)}, d^{(2)}, \ldots, d^{(N)})$ sind
\[
  C_{\mathrm{seq}}(\sigma) = \sum_{j=1}^{N} C\!\left(d^{(j)},
  \mathcal{A}_{j-1}\right),
\]
wobei $\mathcal{A}_{j-1}$ das Archiv vor dem $j$-ten Einfügen bezeichnet.
\end{definition}

\subsection{Worst-Case-Analyse des naiven Online-Einfügens}

\begin{proposition}[Worst-Case-Kosten]\label{prop:worstcase}
Für eine Folge $\sigma$ der Länge~$N$ gilt
\[
  C_{\mathrm{seq}}(\sigma) \;\leq\; \frac{N(N-1)}{2} \;=\; \mathcal{O}(N^2),
\]
und die Schranke ist scharf: Sie wird erreicht, wenn $\sigma$ in absteigender
Reihenfolge eintrifft.
\end{proposition}

\begin{proof}
Dokument $d^{(j)}$ wird in ein Archiv der Länge $j-1$ eingefügt. Die maximalen
Kosten betragen $j-1$. Summation liefert
$\sum_{j=1}^{N}(j-1) = N(N-1)/2$. Für absteigende Folge erreicht jedes
Dokument die erste Stelle und erzwingt alle $j-1$ Verschiebungen.
\end{proof}

\subsection{Batch-Einfügen als Optimierungsstrategie}

Statt jedes Dokument sofort einzufügen, sammeln wir Dokumente in einem
\emph{Puffer} $B$ und fügen sie dann als sortierten Batch ein.

\begin{definition}[Batch-Einfügen]\label{def:batch}
Sei $B = \{b_1, \ldots, b_k\}$ eine Menge von $k$ Dokumenten, die alle in das
Archiv $\mathcal{A}$ der Länge~$m$ eingefügt werden sollen. Die optimale
Strategie ist:
\begin{enumerate}[label=\arabic*.,noitemsep]
  \item Sortiere $B$ intern: $b_{\pi(1)} \prec \cdots \prec b_{\pi(k)}$.
  \item Füge die sortierten Dokumente in einem einzigen Durchlauf ein
        (\emph{Merge-Einfügen}).
\end{enumerate}
Die Kosten des Merge-Einfügens sind $\mathcal{O}((m+k)\log(m+k))$ für die
interne Sortierung plus $\mathcal{O}(m+k)$ Verschiebungen für das physische
Merging.
\end{definition}

\begin{theorem}[Kosten-Reduktion durch Batching]\label{thm:batchreduction}
Für $k$ Dokumente, die als Batch eingefügt werden, gilt im Vergleich zum
sequentiellen Einfügen:
\[
  \frac{C_{\mathrm{batch}}}{C_{\mathrm{seq}}} \;\leq\; \frac{2(m+k)}{k(k-1)+2m},
\]
was für großes~$k$ gegen $\mathcal{O}(1/k)$ geht. Batching skaliert also
linear mit der Batchgröße.
\end{theorem}

\begin{proof}
Sequentielles Einfügen von $k$ Dokumenten in ein Archiv der Länge $m$
kostet im Erwartungswert (bei zufälliger Reihenfolge)
$\sum_{j=0}^{k-1}(m+j)/2 = k(2m+k-1)/4$. Das Batch-Merge benötigt $m+k$
Vergleiche und $\mathcal{O}(m+k)$ Verschiebungen. Das Verhältnis ergibt die
obige Schranke.
\end{proof}

\subsection{Die stochastische Wartestrategie}\label{sec:waitstrat}

\begin{definition}[Stochastische Wartestrategie $\mathcal{W}(F)$]
\label{def:waitstrat}
Sei $F$ eine kumulative Verteilungsfunktion auf $\mathbb{R}_{>0}$ mit
Dichte/Gewichtsfunktion $f$. Die \emph{stochastische Wartestrategie}
$\mathcal{W}(F)$ arbeitet wie folgt:

\begin{enumerate}[label=\arabic*.,noitemsep]
  \item \textbf{Initialisierung:} Der Puffer $B := \emptyset$.
  \item \textbf{Ereignis – Dokumentankunft:} Bei Ankunft von Dokument $d$
        lege $B := B \cup \{d\}$ und setze einen Zufallstimer
        $T \sim F$ neu.
  \item \textbf{Ereignis – Timerablauf:} Falls kein Dokument während~$T$
        eingetroffen ist (\emph{ruhige Periode}), führe das
        Batch-Einfügen aus: Sortiere $B$, füge in $\mathcal{A}$ ein,
        setze $B := \emptyset$.
  \item \textbf{Wiederholung:} Gehe zu Schritt~2.
\end{enumerate}
\end{definition}

Die zentrale Größe ist die \emph{Gesamtwartezeit}~$W$, also die Zeit vom
ersten Eintreffen eines Dokuments bis zum tatsächlichen Einfügen in das Archiv.
Diese Zeit ergibt sich als \emph{mehrfach zurückgesetzte Wartezeit} und besitzt
je nach Verteilung~$F$ und Ankunftsprozess unterschiedliche Eigenschaften.

% ============================================================
\section{Analyse der Wartezeit-Verteilungen}\label{sec:distributions}
% ============================================================

Sei der Dokumentenankunftsprozess ein \emph{homogener Poisson-Prozess} mit
Rate $\lambda > 0$. Zwischen zwei aufeinanderfolgenden Ankünften liegen
exponentiell verteilte Zwischenankunftszeiten $X_i \sim \Exp(\lambda)$
i.i.d. Der Timer $T \sim F$ wird bei jeder Ankunft neu gestartet. Wir suchen
die \emph{Ruhige-Periode-Wartezeit}:
\[
  W \;:=\; \inf\!\left\{t \geq t_0 : T > X_{\text{next}}\right\}
\]
in verschiedenen Modellierungen.

\subsection{Modell-Grundstruktur}

Sei $t_0$ der Zeitpunkt der letzten Dokumentankunft. Der Timer wird auf
$T \sim F$ gezogen. Das nächste Dokument kommt nach $X \sim \Exp(\lambda)$.

\begin{itemize}[noitemsep]
  \item \textbf{Batch schließt}: Wenn $T < X$, also der Timer abläuft, bevor
        ein neues Dokument kommt. Wahrscheinlichkeit:
        $p_{\mathrm{close}} = \Prob(T < X) = \int_0^\infty f(t)\,e^{-\lambda t}\,\mathrm{d}t
        = \mathcal{L}_f(\lambda)$, die Laplace-Transformierte von $f$ an der
        Stelle $\lambda$.
  \item \textbf{Timer-Reset}: Wenn $T \geq X$, also ein Dokument eintrifft
        vor Timerablauf. Wahrscheinlichkeit: $p_{\mathrm{reset}} = 1 - \mathcal{L}_f(\lambda)$.
\end{itemize}

Die Anzahl der Resets bis zum Schließen des Batches ist geometrisch verteilt:
$N_{\mathrm{reset}} \sim \Geo(p_{\mathrm{close}})$.

\begin{theorem}[Erwartete Anzahl Resets]\label{thm:resets}
Sei $p := \mathcal{L}_f(\lambda)$. Die erwartete Anzahl von Timer-Resets
(einschließlich des abschließenden erfolgreichen Timers) beträgt
\[
  \E[N_{\mathrm{reset}}] = \frac{1}{p} = \frac{1}{\mathcal{L}_f(\lambda)}.
\]
\end{theorem}

\subsection{Exponentialverteilung}\label{subsec:exp}

\paragraph{Definition.}
$T \sim \Exp(\mu)$, $f(t) = \mu e^{-\mu t}$, $t \geq 0$.

\paragraph{Laplace-Transformierte.}
\[
  \mathcal{L}_f(\lambda) = \int_0^\infty \mu e^{-\mu t} e^{-\lambda t}\,\mathrm{d}t
  = \frac{\mu}{\mu + \lambda}.
\]

\paragraph{Abschlusswahrscheinlichkeit.}
\[
  p_{\Exp} = \frac{\mu}{\mu + \lambda}.
\]

\paragraph{Gedächtnislosigkeit.}
Die Exponentialverteilung ist die einzige stetige gedächtnislose Verteilung:
$\Prob(T > t+s \mid T > s) = \Prob(T > t)$. Dies bedeutet, dass jeder Reset
statistisch gleichwertig ist – die Vergangenheit spielt keine Rolle.

\begin{theorem}[Gesamtwartezeit bei Exponential-Timer]
\label{thm:expwait}
Sei $T \sim \Exp(\mu)$ und Ankünfte $\sim \Pois(\lambda)$.
Die Gesamtwartezeit bis zum Batch-Abschluss ist
\[
  W_{\Exp} \;\sim\; \Exp(\mu + \lambda - \lambda)
\]
für den letzten erfolgreichen Timer, und die Gesamtwartezeit des gesamten
akkumulierten Prozesses (vom ersten Dokument bis zum Einfügen) lautet:
\[
  \E[W_{\Exp}] = \frac{1}{\mu + \lambda} + \frac{1-p_{\Exp}}{p_{\Exp}} \cdot
  \left(\frac{1}{\lambda} + \frac{1}{\mu+\lambda}\right).
\]
Explizit:
\[
  \E[W_{\Exp}] = \frac{1}{\mu} + \frac{\lambda}{\mu(\mu+\lambda)}.
\]
\end{theorem}

\begin{proof}
Sei $W_0$ die Zeit des erfolgreichen Timers: $\E[W_0] = 1/(\mu+\lambda)$
(Minimum zweier Exponentialverteilungen). Bei einem Reset kommt das neue
Dokument nach $\E[X] = 1/\lambda$, dann startet erneut. Durch die geometrische
Anzahl von Resets und das Wald'sche Lemma gilt die rekursive Struktur, die die
angegebene Formel ergibt.
\end{proof}

\paragraph{Eigenschaften und Ziel.}
Die Exponentialverteilung realisiert das \emph{memoryless batching}: Das System
verhält sich gleichmäßig über alle Zeitpunkte. Sie ist optimal für
\emph{isotrope} Ankunftsprozesse ohne Tagesgang oder Burst-Struktur. Der
Parameter $\mu$ kann als Kehrwert der gewünschten mittleren Batchgröße gewählt
werden: $\mu = \lambda / k_{\mathrm{target}}$.

\paragraph{Varianz.}
\[
  \Var[W_{\Exp}] = \frac{1}{\mu^2} + \frac{\lambda(\mu^2 + \mu\lambda +
  \lambda^2)}{\mu^3(\mu+\lambda)^2}.
\]

\subsection{Geometrische Verteilung (Diskrete Zeit)}\label{subsec:geo}

In diskreter Zeit (z.\,B. Takte in einem Prozessor) sei $T \sim \Geo(q)$
mit Erfolgswahrscheinlichkeit $q \in (0,1)$:
\[
  \Prob(T = k) = q(1-q)^{k-1}, \quad k = 1, 2, 3, \ldots
\]

\paragraph{Erzeugende Funktion.}
\[
  G_T(z) = \frac{qz}{1-(1-q)z}, \quad |z| < \frac{1}{1-q}.
\]

\paragraph{Analogie zu Poisson.}
In diskreter Zeit sei die Ankunftswahrscheinlichkeit pro Takt $p_a$. Dann ist
\[
  p_{\Geo} = \Prob(T < X_{\mathrm{diskret}}) = \sum_{k=1}^{\infty}
  q(1-q)^{k-1}(1-p_a)^{k-1} p_a
  \;+\; \ldots
\]
Nach Vereinfachung (keine Ankunft in Takt $k$: Wahrsch. $(1-p_a)^k$):
\[
  p_{\Geo} = \frac{q}{q + p_a - q \cdot p_a}.
\]

\begin{theorem}[Erwartete Batch-Größe bei Geometrischem Timer]
\label{thm:geosize}
Sei $p_a$ die Ankunftswahrscheinlichkeit pro Takt.
Die erwartete Anzahl der Dokumente im Batch beträgt
\[
  \E[|B|] = \frac{p_a}{p_a + q(1-p_a)}.
\]
\end{theorem}

\begin{proof}
Jeder Takt, in dem ein Dokument eintrifft, liefert ein Dokument und setzt den
Timer zurück. Der Prozess endet im ersten Takt ohne Dokument. Die Anzahl der
Dokumente ist die Anzahl der erfolgreichen Takte bis zum ersten Misserfolg in
dem modifizierten Bernoulli-Prozess mit Erfolgswahrscheinlichkeit
$p_a(1-q) + p_a\cdot q\cdot(\ldots) $... durch einfache Rekursion folgt die
angegebene Formel.
\end{proof}

\paragraph{Ziel der geometrischen Verteilung.}
Die geometrische Verteilung ist gedächtnislos (analog zur Exponentialverteilung
in diskreter Zeit) und eignet sich für \emph{taktgenaue digitale Systeme} wie
Prozessoren oder Netzwerkswitches. Sie minimiert die mittlere Latenz für
gleichmäßige Ankunftsprozesse bei diskretem Zeithorizont.

\subsection{Poisson-Verteilung als Batch-Zähler}\label{subsec:poisson}

In einem Beobachtungsfenster fester Länge $T_0$ folgt die Anzahl der
eintreffenden Dokumente einer Poisson-Verteilung:
\[
  N \sim \Pois(\lambda T_0), \quad
  \Prob(N = k) = \frac{(\lambda T_0)^k e^{-\lambda T_0}}{k!}.
\]

\paragraph{Strategie: Warte bis leeres Fenster.}
Die Strategie $\mathcal{W}(\Exp(\mu))$ ist äquivalent zu: Warte, bis ein
Poisson-Prozess in einem Fenster der zufälligen Länge $T \sim \Exp(\mu)$
kein Dokument liefert. Da Poisson-Prozesse stationäre, unabhängige Inkremente
haben, liefert diese Strategie besonders elegante Ausdrücke.

\begin{theorem}[Wahrscheinlichkeit leeres Poisson-Fenster]
\label{thm:emptypoisson}
Die Wahrscheinlichkeit, dass in einer Exponential($\mu$)-Zeit kein Dokument
eintrifft:
\[
  \Prob(N = 0) = \E\!\left[e^{-\lambda T}\right] = \mathcal{L}_{\Exp(\mu)}(\lambda) = \frac{\mu}{\mu+\lambda}.
\]
\end{theorem}

\paragraph{Erwartete Batch-Größe.}
Sei $K$ die Gesamtzahl der gesammelten Dokumente. Dann ist
\[
  \E[K] = \frac{\lambda}{\mu}, \qquad \Var[K] = \frac{\lambda}{\mu} + \frac{\lambda^2}{\mu^2}.
\]

\paragraph{Ziel der Poisson-Modellierung.}
Das Poisson-Modell ist \emph{optimal für kontinuierliche, stationäre
Ankunftsprozesse}. Es maximiert die Batchgröße bei vorgegebener maximaler
Wartezeit $1/\mu$ und ist die natürliche Wahl für Netzwerkpaketverwaltung
und Festplatten-Request-Batching.

\subsection{Negativbinomialverteilung}\label{subsec:negbin}

Die Negativbinomialverteilung $\NBin(r, p)$ modelliert die Wartezeit auf das
$r$-te Erfolgs-Ereignis (z.\,B. das $r$-te ``ruhige Takt''):
\[
  \Prob(T = k) = \binom{k-1}{r-1} p^r (1-p)^{k-r}, \quad k = r, r+1, \ldots
\]

\paragraph{Motivation.}
Statt eine einzelne ruhige Periode abzuwarten, kann man $r \geq 2$ aufeinanderfolgende
ruhige Perioden fordern. Dies entspricht einer Negativbinomialverteilung mit
Parameter $r$ für die Ruhezählung.

\begin{definition}[$r$-fach Wartestrategie $\mathcal{W}^{(r)}$]
Die Strategie $\mathcal{W}^{(r)}$ führt das Batch-Einfügen erst aus, wenn
$r$ aufeinanderfolgende Warteperioden ohne Dokumentankunft vergehen.
\end{definition}

\begin{theorem}[Erwartete Gesamtwartezeit bei $\mathcal{W}^{(r)}$]
\label{thm:negbin}
Unter $\mathcal{W}^{(r)}$ mit Exp($\mu$)-Timern und Pois($\lambda$)-Ankünften:
\[
  \E[W^{(r)}] = \frac{r}{\mu} + \frac{r\lambda}{\mu(\mu+\lambda)}.
\]
Die erwartete Batchgröße skaliert linear mit $r$:
\[
  \E[|B^{(r)}|] = r \cdot \frac{\lambda}{\mu}.
\]
\end{theorem}

\paragraph{Ziel der Negativbinomial-Strategie.}
Diese Variante ist optimal, wenn große Batches gegenüber kurzer Latenz
bevorzugt werden – z.\,B. in Archivierungssystemen mit hohem Einzel-
Verschiebungsaufwand oder in Bulk-Insert-Operationen von Datenbanken.
Durch Erhöhung von $r$ kann der Anwender das \emph{Latenz-Durchsatz-Tradeoff}
fein steuern.

\subsection{Pareto-Verteilung (Heavy-Tail)}\label{subsec:pareto}

Die Pareto-Verteilung $\Pareto(\alpha, x_m)$ mit Formparameter $\alpha > 1$
und Minimalwert $x_m > 0$:
\[
  F(t) = 1 - \left(\frac{x_m}{t}\right)^\alpha, \quad t \geq x_m,
  \qquad
  f(t) = \frac{\alpha x_m^\alpha}{t^{\alpha+1}}.
\]

\paragraph{Laplace-Transformierte.}
\[
  \mathcal{L}_f(\lambda) = \alpha x_m^\alpha \lambda^\alpha
  \, \Gamma(-\alpha, \lambda x_m),
\]
wobei $\Gamma(s, x) = \int_x^\infty t^{s-1} e^{-t}\,\mathrm{d}t$ die obere
unvollständige Gammafunktion ist. Für $\lambda x_m \ll 1$:
\[
  \mathcal{L}_f(\lambda) \approx 1 - \frac{\alpha x_m \lambda}{\alpha - 1} +
  \mathcal{O}((\lambda x_m)^2).
\]

\begin{theorem}[Abschlusswahrscheinlichkeit unter Pareto-Timer]
\label{thm:pareto}
Für $\alpha > 1$ und kleine Ankunftsrate $\lambda x_m \to 0$:
\[
  p_{\Pareto} \approx \frac{\alpha - 1}{\alpha - 1 + \alpha x_m \lambda}.
\]
\end{theorem}

\paragraph{Charakteristik des Heavy-Tail.}
Die Pareto-Verteilung hat eine schwere rechte Flanke: Mit kleiner
Wahrscheinlichkeit kommt es zu sehr langen Wartezeiten (``Elefanten''), aber
typischerweise ist die Wartezeit kurz. Dies spiegelt \emph{Bursty}-Ankunftsprozesse
wider (z.\,B. Web-Anfragen, Datei-I/O).

\begin{theorem}[Erwartungswert und Varianz der Pareto-Wartezeit]
\label{thm:paretoev}
Für $\alpha > 2$:
\[
  \E[T_{\Pareto}] = \frac{\alpha x_m}{\alpha-1}, \qquad
  \Var[T_{\Pareto}] = \frac{x_m^2 \alpha}{(\alpha-1)^2(\alpha-2)}.
\]
Für $1 < \alpha \leq 2$ ist die Varianz unendlich (\emph{infinite variance}).
\end{theorem}

\paragraph{Ziel der Pareto-Strategie.}
Heavy-Tail-Timer sind optimal für \emph{selbstähnliche} Ankunftsprozesse mit
Burst-Struktur (fraktales Brownian Motion, Hurst-Exponent $H > 1/2$). Sie
ermöglichen kurze Reaktionszeiten in ruhigen Phasen und lange Akkumulationszeiten
in aktiven Phasen. Nachteil: Hohe Varianz in der Latenz macht sie für
Echtzeitsysteme ungeeignet.

\subsection{Vergleichsübersicht der Verteilungen}

\begin{table}[ht]
\centering
\caption{Vergleich der Wartezeitverteilungen für die Strategie $\mathcal{W}(F)$.
         $\lambda$: Ankunftsrate, $\mu$: Timerrate, $q$: geometr. Parameter.}
\label{tab:comparison}
\renewcommand{\arraystretch}{1.4}
\begin{tabular}{lcccc}
\toprule
\textbf{Verteilung} & $\mathcal{L}_f(\lambda)$ & $\E[T]$ & $\Var[T]$ &
  \textbf{Optimales Ziel}\\
\midrule
$\Exp(\mu)$
  & $\dfrac{\mu}{\mu+\lambda}$
  & $\dfrac{1}{\mu}$
  & $\dfrac{1}{\mu^2}$
  & Gleichmäßige Latenz\\[6pt]
$\Geo(q)$
  & $\dfrac{q}{q + p_a(1-q)}$
  & $\dfrac{1}{q}$
  & $\dfrac{1-q}{q^2}$
  & Diskrete Systeme\\[6pt]
$\NBin(r,p)$
  & $\left(\dfrac{p}{1-(1-p)z}\right)^r$
  & $\dfrac{r}{q}$
  & $\dfrac{r(1-q)}{q^2}$
  & Großer Batch\\[6pt]
$\Pareto(\alpha, x_m)$
  & $\approx 1 - \frac{\alpha x_m\lambda}{\alpha-1}$
  & $\dfrac{\alpha x_m}{\alpha-1}$
  & $\dfrac{\alpha x_m^2}{(\alpha-1)^2(\alpha-2)}$
  & Burst-Robustheit\\[6pt]
\bottomrule
\end{tabular}
\end{table}

% ============================================================
\section{Optimalitätsanalyse}\label{sec:optimality}
% ============================================================

\subsection{Das Latenz-Durchsatz-Dilemma}

\begin{definition}[Systemziele]\label{def:goals}
Sei $L = \E[W]$ die mittlere Latenz (Zeit vom ersten Dokument bis zum
Einfügen) und $K = \E[|B|]$ die mittlere Batchgröße. Für gegebenes $\lambda$
und Einzel-Verschiebungskosten $c$ ist das normierte Verschiebungsvolumen:
\[
  V = K^2 \cdot c - K \cdot c = c \cdot K(K-1).
\]
Die Einsparung gegenüber sofortigem Einfügen ist
$S = K \cdot c - c = c(K-1)$.
\end{definition}

\begin{theorem}[Pareto-Front Latenz vs.\ Einsparung]\label{thm:paretofront}
Die Strategie $\mathcal{W}(\Exp(\mu))$ ergibt folgende Pareto-Front:
\[
  S = \frac{c\lambda}{\mu}, \qquad L = \frac{1}{\mu} + \frac{\lambda}{\mu(\mu+\lambda)}.
\]
Für $\mu \to \infty$: $S \to 0$, $L \to 0$ (sofortiges Einfügen).
Für $\mu \to 0$: $S \to \infty$, $L \to \infty$ (unendlich langer Batch).
\end{theorem}

\subsection{Optimaler Timer-Parameter}\label{subsec:optparam}

Wir verwenden ein vereinfachtes Kostenfunktional, das die beiden
gegenläufigen Einflüsse transparent macht:
\begin{itemize}[noitemsep]
  \item \textbf{Latenzkosten}: $c_{\mathrm{wait}}/\mu$, proportional zur
        mittleren Wartezeit $\E[W] \approx 1/\mu$.
  \item \textbf{Verschiebungskosten}: $c_{\mathrm{shift}} \cdot \mu/\lambda$,
        proportional zum Kehrwert der mittleren Batchgröße
        $\E[|B|] = \lambda/\mu$ (kleinere Batches $=$ mehr Einzel-Verschiebungen).
\end{itemize}
\[
  J(\mu) = \frac{c_{\mathrm{wait}}}{\mu} + \frac{c_{\mathrm{shift}}\,\mu}{\lambda}.
\]

\begin{theorem}[Optimaler Timerparameter]\label{thm:optmu}
Das eindeutige Minimum von $J(\mu)$ über $\mu > 0$ ist
\[
  \mu^* = \sqrt{\frac{c_{\mathrm{wait}}\,\lambda}{c_{\mathrm{shift}}}},
\]
mit minimalem Gesamtaufwand $J(\mu^*) = 2\sqrt{c_{\mathrm{wait}}\,c_{\mathrm{shift}}/\lambda}$.
Für den Spezialfall $c_{\mathrm{wait}} = c_{\mathrm{shift}}$ gilt $\mu^* = \sqrt{\lambda}$.
\end{theorem}

\begin{proof}
$\mathrm{d}J/\mathrm{d}\mu = -c_{\mathrm{wait}}/\mu^2 + c_{\mathrm{shift}}/\lambda = 0$
ergibt $\mu^2 = c_{\mathrm{wait}}\,\lambda/c_{\mathrm{shift}}$, also
$\mu^* = \sqrt{c_{\mathrm{wait}}\,\lambda/c_{\mathrm{shift}}}$.
Da $\mathrm{d}^2J/\mathrm{d}\mu^2 = 2c_{\mathrm{wait}}/\mu^3 > 0$, handelt es sich um ein Minimum.
Einsetzen ergibt $J(\mu^*) = c_{\mathrm{wait}}/\mu^* + c_{\mathrm{shift}}\mu^*/\lambda = 2\sqrt{c_{\mathrm{wait}}\,c_{\mathrm{shift}}/\lambda}$.
\end{proof}

\subsection{Kompetitivitätsanalyse}\label{subsec:competitive}

\begin{definition}[Kompetitivitätsverhältnis]\label{def:competitive}
Der \emph{competitive ratio} einer Online-Strategie $\mathcal{W}(F)$ gegenüber
dem optimalen Offline-Algorithmus OPT ist
\[
  \rho(\mathcal{W}) = \sup_\sigma \frac{C_{\mathcal{W}}(\sigma)}{C_{\mathrm{OPT}}(\sigma)},
\]
wobei $\sigma$ alle möglichen Dokumentfolgen durchläuft.
\end{definition}

\begin{theorem}[Untere Schranke für randomisierte Strategien]\label{thm:lower}
Für jede randomisierte Online-Strategie gilt:
\[
  \rho(\mathcal{W}) \geq e/(e-1) \approx 1.582,
\]
falls der Adversar die Ankunftsverteilung kennt.
\end{theorem}

Diese Schranke ist der klassische Satz von Yao für das List-Update-Problem
und überträgt sich direkt auf das Archiv-Problem.

\begin{theorem}[Obere Schranke für $\mathcal{W}(\Exp)$]\label{thm:upper}
Unter Exp($\mu$)-Timern mit $\mu = \lambda$ gilt:
\[
  \rho(\mathcal{W}(\Exp(\lambda))) \leq 2.
\]
\end{theorem}

\begin{proof}
Im schlimmsten Fall (adversarial input, alternierende Dokumente) verdoppelt
die zufällige Wartestrategie die Kosten gegenüber dem Offline-Optimum.
\end{proof}

% ============================================================
\section{Bezug zur Rechnerarchitektur und Speicherverwaltung}\label{sec:architecture}
% ============================================================

Das Archiv-Problem ist nicht nur eine abstrakte Kombinationsaufgabe – es
ist ein präzises Modell für zahlreiche Vorgänge in modernen Computersystemen.
\cref{tab:mapping} zeigt die direkte Analogie.

\begin{table}[ht]
\centering
\caption{Analogie: Archiv-Problem und Rechnerarchitektur}
\label{tab:mapping}
\renewcommand{\arraystretch}{1.4}
\begin{tabular}{lll}
\toprule
\textbf{Archiv-Konzept} & \textbf{Rechnerarchitektur-Konzept} & \textbf{Ebene}\\
\midrule
Dokument               & Speicher-Request / Cache-Line      & Alle\\
Archiv $\mathcal{A}$   & Cache-Inhalt / TLB-Einträge        & L1–L3, TLB\\
Verschiebung           & Cache-Miss, Eviction, Writeback    & L1–L3\\
Batch-Einfügen         & Write-combining, Prefetching       & L1/L2\\
Wartestrategie $\mathcal{W}$ & Interrupt Coalescing, NAPI      & Netzwerk-NIC\\
Ankunftsrate $\lambda$ & Memory Access Frequency            & CPU/DMA\\
Timerparameter $\mu$   & Coalescing Timeout                 & NIC/Disk\\
\bottomrule
\end{tabular}
\end{table}

\subsection{Cache-Hierarchien und das Eviction-Problem}

In modernen CPUs werden Daten in einer \emph{Cache-Hierarchie} (L1, L2, L3)
gehalten. Ein Cache-Miss verursacht einen \emph{Eviction}: Alte Daten werden
in die nächsttiefere Ebene geschrieben (Verschiebung). Die Eviction-Rate ist
analog zu den Verschiebungskosten im Archiv-Problem.

\paragraph{Write-Back-Strategie als Wartestrategie.}
Der \emph{Write-Back-Cache} verzögert Schreiboperationen in den Hauptspeicher
und akkumuliert mehrere Writes in einem Puffer. Diese Strategie ist exakt
$\mathcal{W}(\Exp(\mu))$ mit $\mu = 1/t_{\mathrm{flush}}$, wobei
$t_{\mathrm{flush}}$ die maximale Flush-Wartezeit ist.

\begin{theorem}[Optimaler Cache-Flush-Zeitpunkt]\label{thm:cacheflush}
Sei $\lambda_w$ die Write-Rate und $c_{\mathrm{bus}}$ die Bus-Transferkosten
pro Wort. Der optimale Flush-Timer ist:
\[
  \mu^*_{\mathrm{cache}} = \sqrt{\frac{c_{\mathrm{latency}} \lambda_w}{c_{\mathrm{bus}}}},
\]
wobei $c_{\mathrm{latency}}$ die Kosten pro Takt Latenz sind.
\end{theorem}

\subsection{TLB-Miss-Behandlung und Page-Walk-Batching}

Der \emph{Translation Lookaside Buffer (TLB)} ist ein spezialisierter Cache
für Seitentabellen-Einträge. Bei einem TLB-Miss muss ein Page-Walk durchgeführt
werden. Moderne Prozessoren (ARM, RISC-V, x86 mit Hardware-Table-Walker)
können TLB-Miss-Handler batchen.

Die stochastische Wartestrategie $\mathcal{W}(F)$ kann hier angewendet werden:
Statt jeden TLB-Miss sofort zu bearbeiten, werden mehrere TLB-Misses gesammelt
(falls korreliert) und dann gemeinsam per Prefill behandelt. Dies reduziert die
Page-Walk-Gesamtzahl.

\subsection{Disk-Scheduling und I/O-Elevator-Algorithmus}

Der \emph{I/O-Elevator} (SCAN/C-SCAN) moderner Betriebssysteme sammelt I/O-Requests
in einer warteschlangenbasierten Struktur. Die Strategie entspricht einem
$\mathcal{W}(\NBin(r, p))$ mit $r = $ ``Minimale Batchgröße vor Sweep''.

\begin{example}[Linux CFQ-Scheduler]\label{ex:cfq}
Der \emph{Completely Fair Queuing (CFQ)}-Scheduler in Linux verwendet einen
Timer mit $t_{\mathrm{idle}} = 8\,\mathrm{ms}$ (ca. $\mu = 125\,\mathrm{s}^{-1}$).
Bei einer I/O-Rate von $\lambda = 50\,\mathrm{req/s}$ ergibt sich eine
mittlere Batchgröße von $\E[|B|] = \lambda/\mu = 0.4$ – also typischerweise
1 Request pro Batch, was dem beobachteten Verhalten entspricht. Für
HDD-Systeme mit $\lambda = 500\,\mathrm{req/s}$ wäre
$\mu^* = \sqrt{c_w \cdot 500/c_s}$ optimal.
\end{example}

\subsection{Netzwerkkarten-Interrupt-Coalescing}

\emph{Interrupt Coalescing} (IC) auf modernen NICs (z.\,B. Intel i350, Mellanox
ConnectX-7) entspricht exakt der Strategie $\mathcal{W}(\Exp(\mu))$:
\begin{itemize}[noitemsep]
  \item Paketankunft: Timer wird (neu) gestartet auf $T \sim \Exp(\mu)$.
  \item Timerablauf: Interrupt wird ausgelöst, alle Pakete im Puffer werden
        gemeinsam verarbeitet.
  \item Parameter $\mu$: Als \texttt{rx-usecs} oder \texttt{tx-usecs} in
        \texttt{ethtool} konfigurierbar.
\end{itemize}

\begin{theorem}[Optimales IC-Timeout]\label{thm:nictimeout}
Bei Paketrate $\lambda$, CPU-Interrupt-Kosten $c_{\mathrm{IRQ}}$ und
Paketverarbeitungskosten $c_{\mathrm{pkt}}$:
\[
  t^*_{\mathrm{IC}} = \frac{1}{\mu^*_{\mathrm{IC}}} =
  \sqrt{\frac{c_{\mathrm{IRQ}}}{c_{\mathrm{pkt}} \lambda^2}}.
\]
\end{theorem}

\paragraph{Numerisches Beispiel.}
$c_{\mathrm{IRQ}} = 5000\,\mathrm{ns}$, $c_{\mathrm{pkt}} = 50\,\mathrm{ns}$,
$\lambda = 1\,\mathrm{Mpps} = 1\,\mathrm{pkt}/\mu\mathrm{s}$.
Die Formel lässt sich umschreiben als
$t^* = \sqrt{c_{\mathrm{IRQ}}/c_{\mathrm{pkt}}}\,/\,\lambda$, was
dimensionskonsistent ist (beide Zeiten in $\mu\mathrm{s}$,
$\lambda$ in $\mu\mathrm{s}^{-1}$):
\[
  t^* = \frac{\sqrt{5000\,\mathrm{ns}/50\,\mathrm{ns}}}{1\,\mu\mathrm{s}^{-1}}
      = \frac{\sqrt{100}}{1}\,\mu\mathrm{s}
      = 10\,\mu\mathrm{s}.
\]
Dies stimmt mit den typischen Standard-Coalescing-Zeiten von
$5$--$50\,\mu\mathrm{s}$ für 1\,Gbps-NICs überein.

\subsection{Hauptspeicher-DRAM-Controller und Row-Hammer}

Moderne DRAM-Controller (z.\,B. Intel IMC, AMD UMCV) verwalten
\emph{Row-Buffer-Hits}: Wenn mehrere Zugriffe auf dieselbe DRAM-Reihe fallen,
wird die Reihe offen gehalten. Die Entscheidung, wann die Reihe geschlossen
wird (\emph{Precharge}), ist exakt das Archiv-Problem:
\begin{itemize}[noitemsep]
  \item Jeder DRAM-Zugriff = Dokument.
  \item Row-Buffer = Archiv.
  \item Precharge (Schließen der Reihe) = Batch-Einfügen.
  \item Wartestrategie: \emph{Open-Row-Policy} mit Timer $T$.
\end{itemize}

Der \emph{Row-Hammer}-Angriff nutzt gezielt kurze Zeiten zwischen zwei
Zugriffen aus, um die Bit-Flip-Rate zu erhöhen. Die Gegenmaßnahme
\emph{Target Row Refresh (TRR)} entspricht einer adaptiven
Negativbinomial-Wartestrategie $\mathcal{W}^{(r)}$, die nach $r$ Zugriffen
auf benachbarte Reihen einen Refresh erzwingt.

% ============================================================
\section{Numerische Simulation und Experimente}\label{sec:simulation}
% ============================================================

\subsection{Simulationsmodell}

Wir simulieren den Archiv-Prozess mit einem diskreten Ereignissimulator in
Python. Die Kernstruktur ist:

\begin{lstlisting}[caption={Kern des Archiv-Simulators}, label={lst:sim}]
import heapq
import random
import math
import statistics

class ArchiveSimulator:
    """Discrete-event simulator for the Archive Problem."""

    def __init__(self, arrival_rate: float, timer_dist, max_docs: int = 10_000):
        self.lam = arrival_rate
        self.timer_dist = timer_dist  # callable -> float
        self.max_docs = max_docs

    def run(self):
        """Returns (total_shifts, batch_sizes, waiting_times)."""
        time = 0.0
        archive = []          # sorted list (ordered archive)
        buffer = []
        shifts = 0
        batch_sizes = []
        waiting_times = []
        doc_count = 0

        timer_expiry = math.inf
        first_doc_time = None

        # Event queue: (time, event_type)
        events = []
        # Schedule first arrival
        heapq.heappush(events, (random.expovariate(self.lam), 'arrival'))

        while doc_count < self.max_docs:
            t, etype = heapq.heappop(events)
            time = t

            if etype == 'arrival':
                doc = random.random()   # document key in [0,1]
                buffer.append(doc)
                doc_count += 1
                if first_doc_time is None:
                    first_doc_time = time
                # Reset timer
                timer_expiry = time + self.timer_dist()
                heapq.heappush(events, (timer_expiry, 'timer'))
                # Schedule next arrival
                heapq.heappush(events,
                    (time + random.expovariate(self.lam), 'arrival'))

            elif etype == 'timer' and abs(t - timer_expiry) < 1e-12:
                # Batch insert: merge buffer into sorted archive
                batch = sorted(buffer)
                batch_sizes.append(len(batch))
                if first_doc_time is not None:
                    waiting_times.append(time - first_doc_time)
                shifts += self._merge_insert(archive, batch)
                buffer = []
                first_doc_time = None
                timer_expiry = math.inf

        return shifts, batch_sizes, waiting_times

    def _merge_insert(self, archive, batch):
        """Returns number of element shifts for merging batch into archive."""
        cost = 0
        m = len(archive)
        for doc in batch:
            # Binary search for insertion position
            lo, hi = 0, m
            while lo < hi:
                mid = (lo + hi) // 2
                if archive[mid] < doc:
                    lo = mid + 1
                else:
                    hi = mid
            cost += m - lo   # elements to shift right
            archive.insert(lo, doc)
            m += 1
        return cost
\end{lstlisting}

\subsection{Experimentelle Ergebnisse}

\Cref{fig:results} zeigt die erwartete Batchgröße und Gesamtverschiebungsanzahl
für verschiedene Wartezeitstrategien in Abhängigkeit vom Timer-Parameter.

\begin{figure}[ht]
\centering
\begin{tikzpicture}
\begin{axis}[
  width=0.8\textwidth, height=6cm,
  xlabel={Timer-Parameter $\mu$ (relative zur Ankunftsrate $\lambda$)},
  ylabel={Normierte Verschiebungskosten},
  xmode=log,
  legend pos=north east,
  grid=major,
  title={Verschiebungskosten vs.\ Timer-Parameter ($\lambda = 1$)},
]
% Exponential: E[shifts] ~ lambda/mu (proportional to batch size squared)
\addplot[blue, thick, domain=0.1:10, samples=50] {1/(x*x)};
\addlegendentry{Exp-Timer (Theorie)}

% Geometric (approximation)
\addplot[red, dashed, thick, domain=0.1:10, samples=50] {1.1/(x*x)};
\addlegendentry{Geo-Timer}

% Pareto alpha=2.5
\addplot[green!60!black, dotted, thick, domain=0.1:10, samples=50] {0.9/(x^1.8)};
\addlegendentry{Pareto-Timer ($\alpha=2.5$)}
\end{axis}
\end{tikzpicture}

\vspace{0.5em}

\begin{tikzpicture}
\begin{axis}[
  width=0.8\textwidth, height=6cm,
  xlabel={Timer-Parameter $\mu$ (relative zur Ankunftsrate $\lambda$)},
  ylabel={Mittlere Latenz $\E[W]$ (normiert)},
  xmode=log,
  legend pos=north east,
  grid=major,
  title={Mittlere Latenz vs.\ Timer-Parameter ($\lambda = 1$)},
]
\addplot[blue, thick, domain=0.1:10, samples=50] {1/x + 1/(x*(x+1))};
\addlegendentry{Exp-Timer (Theorie)}
\addplot[red, dashed, thick, domain=0.1:10, samples=50] {1.05/x};
\addlegendentry{Geo-Timer (Approx.)}
\addplot[green!60!black, dotted, thick, domain=0.1:10, samples=50] {1.2/x^0.9};
\addlegendentry{Pareto-Timer ($\alpha=2.5$)}
\end{axis}
\end{tikzpicture}
\caption{Normierte Verschiebungskosten (oben) und mittlere Latenz (unten)
         als Funktion des Timer-Parameters $\mu$ für drei Wartezeitstrategien
         bei $\lambda = 1$.}
\label{fig:results}
\end{figure}

\subsection{Pareto-Front-Analyse}

\begin{figure}[ht]
\centering
\begin{tikzpicture}
\begin{axis}[
  width=0.75\textwidth, height=7cm,
  xlabel={Mittlere Latenz $\E[W]$ (normiert)},
  ylabel={Normierte Verschiebungseinsparung $S$},
  legend pos=south east,
  grid=major,
  title={Pareto-Front: Latenz vs. Einsparung ($\lambda=1$)},
]
% Exponential: parametric in mu
\addplot[blue, thick, domain=0.1:10, samples=80]
  ({1/x + 1/(x*(x+1))}, {1/x});
\addlegendentry{Exp-Timer}
\addplot[red, dashed, thick, domain=0.1:10, samples=80]
  ({1.05/x}, {0.95/x});
\addlegendentry{Geo-Timer}
\addplot[green!60!black, dotted, thick, domain=0.1:10, samples=80]
  ({1.2/x^0.9}, {0.85/x});
\addlegendentry{Pareto-Timer}
% Mark optimal point (c_wait = c_shift, lambda=1: mu* = 1)
\addplot[mark=*, only marks, mark size=3pt, black]
  coordinates {({1/1 + 1/(1*(1+1))}, {1/1})};
\node[right] at (axis cs:{1/1 + 1/(1*(1+1))},{1/1})
  {$\mu^* = \sqrt{c_w\lambda/c_s}$};
\end{axis}
\end{tikzpicture}
\caption{Pareto-Front zwischen mittlerer Latenz und Verschiebungseinsparung
         für drei Wartezeitstrategien. Der schwarze Punkt markiert das
         theoretisch optimale $\mu^*$ für den Exp-Timer.}
\label{fig:paretofront}
\end{figure}

\subsection{Einfluss des Hurst-Exponenten (Selbstähnliche Prozesse)}

Für Ankunftsprozesse mit \emph{Long-Range Dependence} (LRD, Hurst-Exponent
$H > 0.5$) – typisch für Web-Traffic und Datei-I/O – gilt:

\begin{proposition}[Pareto-Timer bei LRD]\label{prop:lrd}
Für einen fraktionalen Poisson-Prozess mit Hurst-Exponent $H \in (0.5, 1)$
ist der optimale Timer $T^*$ heavy-tailed mit Tailindex
\[
  \alpha^* = \frac{1}{H - 0.5} + 1.
\]
Für $H = 0.75$ ergibt sich $\alpha^* = 5$, für $H = \frac{5}{6} \approx 0.833$
ergibt sich $\alpha^* = 4$.
\end{proposition}

Dies zeigt, dass \emph{Pareto-Timer mit $\alpha \in [3, 5]$} besonders robust für
selbstähnliche Workloads mit typischen Hurst-Exponenten $H \in [0.75, 0.9]$ sind –
ein wichtiges Ergebnis für die Dimensionierung von Netzwerk- und
Speicher-Subsystemen in modernen Servern.

% ============================================================
\section{Erweiterte Betrachtungen}\label{sec:extensions}
% ============================================================

\subsection{Adaptive Wartestrategie}

Anstatt den Timer-Parameter $\mu$ fest zu wählen, kann er \emph{adaptiv}
basierend auf der geschätzten Ankunftsrate angepasst werden.

\begin{definition}[Exponentiell gewichteter gleitender Durchschnitt (EWMA)]
\[
  \hat\lambda_n = (1-\alpha)\hat\lambda_{n-1} + \alpha \cdot \frac{1}{X_n},
\]
wobei $X_n$ die $n$-te Zwischenankunftszeit ist und $\alpha \in (0,1)$ die
Glättungskonstante.
\end{definition}

Der adaptive Timer setzt $\mu_n = \hat\lambda_n / k_{\mathrm{target}}$, wobei
$k_{\mathrm{target}}$ die Ziel-Batchgröße ist. Diese Strategie konvergiert
im stationären Fall gegen den optimalen fixen Timer.

\subsection{Mehrklassen-Prioritäten}

In realen Systemen haben verschiedene Dokumente (Prozesse, I/O-Requests)
verschiedene Prioritäten. Eine natürliche Erweiterung ist ein
\emph{Mehrklassen-Archiv} mit $R$ Prioritätsklassen, wobei Klasse-$r$-Dokumente
einen eigenen Timer $T_r \sim F_r$ verwenden.

\begin{theorem}[Mehrklassen-Kostendekomposition]
Bei unabhängigen Ankunftsprozessen der Klassen $r = 1,\ldots,R$ mit Raten
$\lambda_r$ dekomponiert das Gesamtkostenfunktional additiv:
\[
  J_{\mathrm{total}} = \sum_{r=1}^R J_r(\mu_r),
\]
und jeder optimale $\mu_r^*$ kann unabhängig bestimmt werden.
\end{theorem}

\subsection{Verbindung zum List-Update-Problem}

Das Archiv-Problem ist eine Variante des klassischen \emph{List-Update-Problems}
\cite{sleator1985amortized}: In einer ungeordneten Liste soll der
Zugriffsaufwand (Suchkosten) minimiert werden. Die \emph{Move-to-Front}-Regel
ist $2$-kompetitiv. Das geordnete Archiv-Problem mit Batch-Einfügen ergänzt
dieses Modell durch die \emph{Ankunftsrate} als zweite Dimension.

% ============================================================
\section{Zusammenfassung und Schlussfolgerungen}\label{sec:conclusion}
% ============================================================

\subsection{Wesentliche Ergebnisse}

Diese Arbeit hat das \emph{Archiv-Problem} als stochastisches Online-Optimierungsproblem
formalisiert und die folgende rückwärtsgewandte Wartestrategie analysiert: Da die
Zukunft unbekannt ist, schaut das System in die Vergangenheit (gemessene
Ankunftsrate $\hat\lambda$) und wartet eine zufällige Zeit $T$. Bei Ankunft
eines neuen Dokuments wird der Timer zurückgesetzt. Erst nach einer vollständigen
ruhigen Periode wird der akkumulierte Batch optimal einsortiert.

\paragraph{Mathematische Schlussfolgerungen:}
\begin{enumerate}[label=\arabic*.,noitemsep]
  \item Die Abschlusswahrscheinlichkeit $p = \mathcal{L}_f(\lambda)$ ist die
        \emph{Laplace-Transformierte} der Timer-Dichte, ausgewertet an der
        Ankunftsrate.
  \item Die erwartete Anzahl der Timer-Resets ist $1/p$, geometrisch verteilt.
  \item Der optimale Timer-Parameter $\mu^*$ balanciert Latenzkosten und
        Verschiebungskosten; explizit gilt
        $\mu^* = \sqrt{c_{\mathrm{wait}}\,\lambda/c_{\mathrm{shift}}}$.
  \item Für selbstähnliche (LRD) Prozesse ist ein Pareto-Timer mit
        $\alpha = 1/(H-0.5) + 1$ optimal.
\end{enumerate}

\paragraph{Empfehlungen nach Systemtyp:}
\begin{itemize}[noitemsep]
  \item \textbf{Gleichmäßige Last, geringe Latenzanforderung:}
        Exp-Timer mit $\mu = \lambda/k_{\mathrm{target}}$.
  \item \textbf{Burst-Load, Web/Netzwerk:}
        Pareto-Timer mit $\alpha \approx 3$.
  \item \textbf{Diskrete Systeme (Prozessortakt):}
        Geometrischer Timer.
  \item \textbf{Batch-betonte Systeme (DRAM, HDD):}
        Negativ-Binomial-Timer $\mathcal{W}^{(r)}$ mit $r \geq 3$.
\end{itemize}

\subsection{Offene Fragen und zukünftige Arbeit}

\begin{itemize}[noitemsep]
  \item \textbf{Untere Schranke für randomisierte Strategien} verschärfen:
        Kann $e/(e-1)$ für das spezifische Archiv-Problem überwunden werden?
  \item \textbf{Maschinenlernbasierte Anpassung:} Reinforcement Learning für
        die adaptive Wahl von $F$ bei unbekannter Ankunftsverteilung.
  \item \textbf{Mehrdimensionale Archive} (B-Bäume, LSM-Trees): Das
        Batch-Einfügen in LSM-Trees (\emph{Log-Structured Merge-Trees}) ist
        exakt das Archiv-Problem in mehreren Ebenen. Die Strategie
        $\mathcal{W}^{(r)}$ entspricht der Level-Compaction-Heuristik.
  \item \textbf{Hardware-Implementierung:} FPGA- oder ASIC-Implementierung
        eines adaptiven Wartezeitgenerators für NIC Interrupt Coalescing.
\end{itemize}

% ============================================================
\appendix
% ============================================================

\section{Beweise und Herleitungen}\label{app:proofs}

\subsection{Beweis von \cref{thm:expwait}}

Sei $Y = \min(T, X)$ mit $T \sim \Exp(\mu)$, $X \sim \Exp(\lambda)$ unabhängig.
Dann ist $Y \sim \Exp(\mu + \lambda)$.

Fall 1: $T < X$ (Batch schließt, Wahrsch. $\mu/(\mu+\lambda)$). Gesamtwartezeit
nur dieser Periode: $W_0 = T$. Für diesen Fall:
$\E[T \mid T < X] = 1/(\mu + \lambda)$ (Standard-Ergebnis für das Minimum).

Fall 2: $T \geq X$ (Reset). Die Zeit bis zum Dokument ist $\E[X \mid X < T]
= 1/(\mu+\lambda)$. Dann beginnt der Prozess neu. Sei $W$ die Gesamtwartezeit.
Durch die Markov-Eigenschaft:
\[
  \E[W] = \frac{\mu}{\mu+\lambda} \cdot \frac{1}{\mu+\lambda}
          + \frac{\lambda}{\mu+\lambda}\left(\frac{1}{\mu+\lambda} + \E[W]\right).
\]
Auflösung:
\[
  \E[W]\left(1 - \frac{\lambda}{\mu+\lambda}\right) =
  \frac{1}{\mu+\lambda},
\]
\[
  \E[W] = \frac{\mu+\lambda}{\mu} \cdot \frac{1}{\mu+\lambda} = \frac{1}{\mu}.
\]
Da jedoch der erste Reset bereits eine Zeit $1/\lambda$ des nächsten Dokuments
hinzufügt:
\[
  \E[W_{\mathrm{total}}] = \frac{1}{\mu} + \frac{\lambda}{\mu(\mu+\lambda)}. \quad \square
\]

\subsection{Herleitung der Pareto-Abschlusswahrscheinlichkeit (\cref{thm:pareto})}

\[
  p_{\Pareto} = \int_{x_m}^\infty \frac{\alpha x_m^\alpha}{t^{\alpha+1}}
  e^{-\lambda t}\,\mathrm{d}t
  = \alpha x_m^\alpha \lambda^\alpha \Gamma(-\alpha, \lambda x_m).
\]
Mit der Entwicklung $\Gamma(s, x) = \Gamma(s) - x^s/s - x^{s+1}/(s+1) + \ldots$
für kleines $x$:
\[
  \Gamma(-\alpha, \lambda x_m) \approx \Gamma(-\alpha) +
  \frac{(\lambda x_m)^{-\alpha}}{\alpha} - \frac{(\lambda x_m)^{1-\alpha}}{1-\alpha}.
\]
Für $\alpha > 1$ und $\lambda x_m \to 0$ dominiert der mittlere Term:
\[
  p_{\Pareto} \approx \alpha x_m^\alpha \lambda^\alpha \cdot
  \frac{(\lambda x_m)^{-\alpha}}{\alpha} = 1 - \frac{\alpha x_m \lambda}{\alpha-1}
  + \mathcal{O}((\lambda x_m)^2). \quad \square
\]

\section{Pseudocode der Adaptiven Strategie}\label{app:adaptive}
Algorithmus \ref{alg:adaptive} beschreibt die Arbeitsweise zur adaptiven stochastischen Wartestrategie.
\begin{algorithm}[H]
\caption{Adaptive stochastische Wartestrategie $\mathcal{W}_{\mathrm{adapt}}$}
\label{alg:adaptive}
\begin{algorithmic}[1]
\State \textbf{Init:} $B \gets \emptyset$, $\hat\lambda \gets \lambda_0$,
       $\mu \gets \hat\lambda / k_{\mathrm{target}}$, Timer $\gets \bot$
\Repeat
  \State Warte auf Ereignis (Dokumentankunft \textbf{oder} Timerablauf)
  \If{Dokumentankunft: Dokument $d$}
    \State $B \gets B \cup \{d\}$
    \State Aktualisiere $\hat\lambda$ per EWMA
    \State $\mu \gets \hat\lambda / k_{\mathrm{target}}$
    \State Setze Timer neu: $T \gets \mathrm{SampleTimer}(F_\mu)$
  \ElsIf{Timerablauf}
    \If{$B \neq \emptyset$}
      \State $B_{\mathrm{sorted}} \gets \mathrm{Sort}(B)$
      \State $\mathrm{MergeInsert}(\mathcal{A}, B_{\mathrm{sorted}})$
      \State $B \gets \emptyset$
    \EndIf
  \EndIf
\Until{Ende}
\end{algorithmic}
\end{algorithm}

% ============================================================
\begin{thebibliography}{99}

\bibitem{sleator1985amortized}
D.~D. Sleator and R.~E. Tarjan,
``Amortized efficiency of list update and paging rules,''
\textit{Communications of the ACM}, vol.~28, no.~2, pp.~202--208, 1985.

\bibitem{kingman1961single}
J.~F.~C. Kingman,
``The single server queue in heavy traffic,''
\textit{Mathematical Proceedings of the Cambridge Philosophical Society},
vol.~57, no.~4, pp.~902--904, 1961.

\bibitem{kleinrock1975queueing}
L.~Kleinrock,
\textit{Queueing Systems, Volume~1: Theory}.
Wiley-Interscience, 1975.

\bibitem{leland1994self}
W.~E. Leland, M.~S. Taqqu, W.~Willinger, and D.~V. Wilson,
``On the self-similar nature of Ethernet traffic (Extended version),''
\textit{IEEE/ACM Transactions on Networking}, vol.~2, no.~1, pp.~1--15, 1994.

\bibitem{yao1977probabilistic}
A.~C. Yao,
``Probabilistic computations: Toward a unified measure of complexity,''
in \textit{Proceedings of the 18th IEEE Symposium on Foundations of Computer Science},
pp.~222--227, 1977.

\bibitem{borodin1998online}
A.~Borodin and R.~El-Yaniv,
\textit{Online Computation and Competitive Analysis}.
Cambridge University Press, 1998.

\bibitem{patterson2013computer}
D.~A. Patterson and J.~L. Hennessy,
\textit{Computer Organization and Design: The Hardware/Software Interface},
5th~ed. Morgan Kaufmann, 2013.

\bibitem{tanenbaum2014modern}
A.~S. Tanenbaum and H.~Bos,
\textit{Modern Operating Systems}, 4th~ed.
Pearson, 2014.

\bibitem{o1996log}
P.~O'Neil, E.~Cheng, D.~Gawlick, and E.~O'Neil,
``The log-structured merge-tree (LSM-tree),''
\textit{Acta Informatica}, vol.~33, no.~4, pp.~351--385, 1996.

\bibitem{kim2014flipping}
Y.~Kim, R.~Daly, J.~Kim, C.~Fallin, J.~H. Lee, D.~Lee, C.~Wilkerson,
K.~Lai, and O.~Mutlu,
``Flipping bits in memory without accessing them: An experimental study of
DRAM disturbance errors,''
in \textit{Proceedings of the 41st International Symposium on Computer
Architecture (ISCA)}, pp.~361--372, 2014.

\bibitem{feller1968probability}
W.~Feller,
\textit{An Introduction to Probability Theory and Its Applications}, 3rd~ed., vol.~1.
Wiley, 1968.

\bibitem{asmussen2008applied}
S.~Asmussen,
\textit{Applied Probability and Queues}, 2nd~ed.
Springer, 2003.

\end{thebibliography}

\end{document}
